\documentclass[11pt,a4paper]{article}

% ===================== POLICES & LANGUE =====================
\usepackage{fontspec}
\setmainfont{Latin Modern Roman}
\setsansfont{Latin Modern Sans}
\setmonofont{Latin Modern Mono}
\usepackage{polyglossia}
\setdefaultlanguage{french}

% ===================== MISE EN PAGE =====================
\usepackage[margin=2.5cm]{geometry}
\usepackage{amsmath,amssymb,graphicx,xcolor,booktabs,hyperref,tcolorbox,enumitem,fancyhdr,array,pifont,longtable,multirow,listings}
\tcbuselibrary{theorems,breakable,skins}
\usepackage{tikz}
\usetikzlibrary{shapes.geometric,arrows.meta,positioning,calc,backgrounds,fit}

\newcommand{\cmark}{\ding{51}}
\newcommand{\xmark}{\ding{55}}
\hypersetup{colorlinks=true,linkcolor=blue!60!black,urlcolor=blue!60!black,citecolor=blue!60!black}

% ===================== BOÎTES =====================
\newtcolorbox{objectifbox}[1][]{colback=blue!5,colframe=blue!50,fonttitle=\bfseries,title={#1},breakable}
\newtcolorbox{infobox}[1][]{colback=green!5,colframe=green!50!black,fonttitle=\bfseries,title={#1},breakable}
\newtcolorbox{attentionbox}[1][]{colback=orange!5,colframe=orange!70,fonttitle=\bfseries,title={#1},breakable}
\newtcolorbox{retenirbox}[1][]{colback=yellow!10,colframe=yellow!50!black,fonttitle=\bfseries,title={#1},breakable}
\newtcolorbox{prereqbox}[1][]{colback=cyan!5,colframe=cyan!50!black,fonttitle=\bfseries,title={#1},breakable}
\newtcolorbox{projetbox}[1][]{colback=teal!5,colframe=teal!50!black,fonttitle=\bfseries,title={#1},breakable}

% ===================== CODE (listings) =====================
\definecolor{codebg}{gray}{0.96}
\definecolor{kwcolor}{RGB}{0,90,160}
\definecolor{strcolor}{RGB}{160,55,30}
\definecolor{comcolor}{RGB}{95,120,95}
\definecolor{ndcolor}{RGB}{120,80,150}

\lstdefinelanguage{JS}{
  morekeywords={let,const,var,function,return,if,else,for,of,in,while,do,true,false,null,undefined,typeof,break,continue,new,this,async,await,try,catch,throw,finally},
  morekeywords=[2]{console,log,table,document,Math,window,alert,prompt,forEach,map,filter,find,push,pop,querySelector,querySelectorAll,addEventListener,classList,textContent,innerHTML,createElement,appendChild,getElementById,includes,fetch,json,JSON,then,stringify,parse,setItem,getItem,removeItem,setInterval,clearInterval,dataset,values,keys,style},
  sensitive=true,
  morecomment=[l]{//},
  morecomment=[s]{/*}{*/},
  morestring=[b]",
  morestring=[b]',
  morestring=[b]`,
}

\lstdefinestyle{js}{
  language=JS,
  basicstyle=\ttfamily\small,
  keywordstyle=\color{kwcolor},
  keywordstyle=[2]\color{ndcolor},
  commentstyle=\color{comcolor}\itshape,
  stringstyle=\color{strcolor},
  backgroundcolor=\color{codebg},
  frame=single,
  rulecolor=\color{gray!45},
  breaklines=true,
  breakatwhitespace=false,
  showstringspaces=false,
  numbers=none,
  tabsize=2,
  columns=fullflexible,
  keepspaces=true,
  literate=%
    {é}{{é}}1 {è}{{è}}1 {à}{{à}}1 {â}{{â}}1 {ê}{{ê}}1
    {î}{{î}}1 {ô}{{ô}}1 {û}{{û}}1 {ù}{{ù}}1 {ç}{{ç}}1
    {ï}{{ï}}1 {ë}{{ë}}1 {É}{{É}}1 {È}{{È}}1 {À}{{À}}1
    {œ}{{œ}}1 {°}{{°}}1 {²}{{²}}1
}

\lstdefinestyle{html}{
  basicstyle=\ttfamily\small,
  backgroundcolor=\color{codebg},
  frame=single,
  rulecolor=\color{gray!45},
  breaklines=true,
  showstringspaces=false,
  numbers=none,
  tabsize=2,
  columns=fullflexible,
  keepspaces=true,
  literate=%
    {é}{{é}}1 {è}{{è}}1 {à}{{à}}1 {ê}{{ê}}1 {ç}{{ç}}1 {î}{{î}}1
}

\lstset{style=js}
\newcommand{\code}[1]{\texttt{#1}}
\setlength{\emergencystretch}{3em}

% ===================== EN-TÊTES =====================
\pagestyle{fancy}
\fancyhf{}
\fancyhead[L]{\small Jàng Afrig --- Lab 9 : Les modes de jeu}
\fancyhead[R]{\small UCAD --- L1 DUT1-Info}
\fancyfoot[C]{\thepage}
\fancyfoot[L]{\small Dr. El Hadji Bassirou TOURÉ}
\fancyfoot[R]{\small 2025--2026}
\renewcommand{\headrulewidth}{0.4pt}
\renewcommand{\footrulewidth}{0.4pt}

% ===================== TITRE =====================
\title{%
\vspace{-1cm}
\textbf{\LARGE Fondamentaux du Web}\\
\textbf{\LARGE HTML, CSS, JavaScript}\\[0.8cm]
\textbf{\huge Lab 9 --- Les modes de jeu}\\[0.3cm]
\textbf{\Large Drapeaux, Langues, Voisins}\\[0.5cm]
\large Licence 1 DUT1-Informatique --- 2025--2026}
\date{}
\author{}

\begin{document}
\maketitle
\thispagestyle{empty}

\begin{prereqbox}[Informations pratiques]
\textbf{Où sommes-nous ?} Jàng Afrig charge les vrais pays et propose un mode Capitales. Mais nos données contiennent bien plus : drapeaux, langues, voisins ! Aujourd'hui, on transforme le quiz en un jeu à \textbf{quatre modes}, sans réécrire le moteur --- juste en variant les questions.

\smallskip
\textbf{Avant de commencer :}
\begin{itemize}[leftmargin=*,noitemsep]
  \item Le Lab 8 fonctionne (tag \code{v0.8}) : les pays se chargent depuis l'API.
  \item Vous avez dans chaque pays : \code{nom}, \code{capitale}, \code{drapeau}, \code{langues}, \code{voisins}.
\end{itemize}

\smallskip
\textbf{Validation de fin :} les quatre modes (Capitales, Drapeaux, Langues, Voisins) fonctionnent, avec une barre de progression. Tag \code{v0.9} poussé + rapport.
\end{prereqbox}

\begin{objectifbox}[Objectifs du Lab 9]
À la fin, vous saurez :
\begin{enumerate}[leftmargin=*,noitemsep]
  \item Proposer plusieurs modes de jeu depuis l'écran d'accueil.
  \item \textbf{Paramétrer} la génération d'une question selon le mode choisi.
  \item Afficher une \textbf{image} (le drapeau) dans une question.
  \item Construire des questions sur les langues et le nombre de voisins.
  \item Ajouter une \textbf{barre de progression} visuelle.
\end{enumerate}
\end{objectifbox}

\tableofcontents
\bigskip\hrule\bigskip

\section*{Rappel express du Lab 8}
\addcontentsline{toc}{section}{Rappel express du Lab 8}
Chaque pays est maintenant un objet riche : \code{\{ cca3, nom, capitale, drapeau, langues, voisins \}}. On va exploiter ces nouveaux champs.

% ===================== SECTION 1 =====================
\section{Choisir un mode depuis l'accueil}

On remplace l'unique bouton « Mode Capitales » par \textbf{quatre boutons}. Chaque bouton porte un attribut \code{data-mode} qui dit lequel c'est. Dans l'écran \code{\#accueil} de \code{index.html} :

\begin{lstlisting}[style=html]
<section id="accueil" class="cachee">
  <p>Choisis ton mode de jeu :</p>
  <label for="pseudo">Ton pseudo :</label>
  <input type="text" id="pseudo" placeholder="Ex : Aminata">

  <div id="modes">
    <button class="mode" data-mode="capitale">Capitales</button>
    <button class="mode" data-mode="drapeau">Drapeaux</button>
    <button class="mode" data-mode="langue">Langues</button>
    <button class="mode" data-mode="voisin">Voisins</button>
  </div>
</section>
\end{lstlisting}

Dans \code{script.js}, on les écoute tous, et on lit le mode dans \code{event.target.dataset.mode} :

\begin{lstlisting}[style=js]
let modeActuel = "capitale";
const boutonsMode = document.querySelectorAll(".mode");

boutonsMode.forEach((bouton) => {
  bouton.addEventListener("click", (event) => {
    const mode = event.target.dataset.mode; // lit l'attribut data-mode
    demarrerQuiz(mode);
  });
});
\end{lstlisting}

\begin{infobox}[Une nouveauté à prendre telle quelle : \texttt{data-mode}]
Un attribut HTML qui commence par \code{data-} sert à \textbf{ranger une info} sur un élément. On le relit en JavaScript avec \code{.dataset} : \code{data-mode} devient \code{element.dataset.mode}. C'est la façon propre d'attacher « ce bouton lance le mode drapeau » directement dans le HTML.
\end{infobox}

On adapte \code{demarrerQuiz} pour qu'il \textbf{reçoive le mode} :

\begin{lstlisting}[style=js]
function demarrerQuiz(mode) {
  modeActuel = mode;
  pseudo = champPseudo.value;
  if (pseudo === "") { pseudo = "Joueur"; }
  score = 0;
  indexQuestion = 0;
  questions = genererQuestions(10, mode);
  afficherEcran(ecranQuiz);
  afficherQuestion();
}

function genererQuestions(nombre, mode) {
  const liste = [];
  for (let i = 0; i < nombre; i++) {
    liste.push(genererQuestion(mode));
  }
  return liste;
}
\end{lstlisting}

% ===================== SECTION 2 =====================
\section{Une question selon le mode}

Le cœur du Lab : une seule fonction \code{genererQuestion(mode)} qui fabrique la bonne question. D'abord, un \textbf{petit outil} qui construit quatre propositions de texte (la bonne + trois distracteurs pris dans un champ donné) :

\begin{lstlisting}[style=js]
// 4 propositions : la bonne réponse + 3 autres prises dans `champ`
function quatrePropositions(bonneReponse, champ) {
  const props = [bonneReponse];
  while (props.length < 4) {
    const autre = choisirAuHasard(pays);
    const valeur = autre[champ]; // ex. autre.capitale ou autre.nom
    if (valeur && !props.includes(valeur)) {
      props.push(valeur);
    }
  }
  return melanger(props);
}
\end{lstlisting}

\code{autre[champ]} utilise l'\textbf{accès par crochets} (Lab 6) : si \code{champ} vaut \code{"nom"}, alors \code{autre[champ]} est \code{autre.nom}. Cela rend l'outil réutilisable pour tous les modes texte.

\subsection{Mode Capitales (rappel) et Drapeaux}
\begin{lstlisting}[style=js]
function genererQuestion(mode) {
  const bonPays = choisirAuHasard(pays);

  if (mode === "capitale") {
    return {
      enonce: `Quelle est la capitale du pays : ${bonPays.nom} ?`,
      bonneReponse: bonPays.capitale,
      propositions: quatrePropositions(bonPays.capitale, "capitale"),
      image: null
    };
  }

  if (mode === "drapeau") {
    return {
      enonce: "À quel pays appartient ce drapeau ?",
      bonneReponse: bonPays.nom,
      propositions: quatrePropositions(bonPays.nom, "nom"),
      image: bonPays.drapeau // l'adresse de l'image du drapeau
    };
  }
  // ... (langue et voisin : ci-dessous)
}
\end{lstlisting}

\subsection{Mode Langues}
On choisit un pays qui a au moins une langue, on en tire une, et on demande qui la parle :

\begin{lstlisting}[style=js]
  if (mode === "langue") {
    let p = choisirAuHasard(pays);
    while (!p.langues) { p = choisirAuHasard(pays); } // s'assurer qu'il a des langues
    const listeLangues = Object.values(p.langues); // ex. ["French", "Wolof"]
    const langue = choisirAuHasard(listeLangues);
    return {
      enonce: `Dans quel pays parle-t-on le ${langue} ?`,
      bonneReponse: p.nom,
      propositions: quatrePropositions(p.nom, "nom"),
      image: null
    };
  }
\end{lstlisting}

\begin{infobox}[\texttt{Object.values}, à prendre tel quel]
Notre champ \code{langues} est un objet comme \code{\{ fra: "French", wol: "Wolof" \}}. \code{Object.values(...)} en extrait la \textbf{liste des valeurs} : \code{["French", "Wolof"]}. On peut alors en choisir une au hasard.
\end{infobox}

\begin{attentionbox}[Une simplification assumée]
Une langue peut être parlée dans plusieurs pays (le français, par exemple). Nos distracteurs étant tirés au hasard, l'un d'eux pourrait, rarement, parler aussi cette langue. À notre niveau, on l'accepte : l'important est de faire fonctionner le mode. On affinera dans une vraie version.
\end{attentionbox}

\subsection{Mode Voisins (réponse numérique)}
Ici la bonne réponse est un \textbf{nombre} (le nombre de voisins). On le transforme en texte, et on fabrique des distracteurs numériques proches :

\begin{lstlisting}[style=js]
  if (mode === "voisin") {
    const nb = bonPays.voisins ? bonPays.voisins.length : 0;
    return {
      enonce: `Combien de pays voisins a ${bonPays.nom} ?`,
      bonneReponse: String(nb),
      propositions: quatreNombres(nb),
      image: null
    };
  }
}

// 4 nombres : le bon + 3 voisins numériques (jamais négatifs)
function quatreNombres(bonNombre) {
  const props = [String(bonNombre)];
  while (props.length < 4) {
    const variation = bonNombre + Math.floor(Math.random() * 7) - 3; // +/- 3
    if (variation >= 0 && !props.includes(String(variation))) {
      props.push(String(variation));
    }
  }
  return melanger(props);
}
\end{lstlisting}

\code{String(nb)} convertit le nombre en texte, car nos boutons affichent et comparent du texte (Lab 7).

% ===================== SECTION 3 =====================
\section{Afficher l'image quand il faut}

Le mode Drapeaux a besoin d'une image. Ajoutez-la dans l'écran \code{\#quiz} de \code{index.html} (cachée par défaut) :

\begin{lstlisting}[style=html]
<section id="quiz" class="cachee">
  <div id="piste-progression"><div id="barre-progression"></div></div>
  <p id="progression">Question 1 / 10</p>
  <img id="image-question" class="cachee" alt="Drapeau à identifier">
  <h2 id="enonce">Énoncé</h2>
  <div id="reponses">
    <button class="reponse">A</button>
    <button class="reponse">B</button>
    <button class="reponse">C</button>
    <button class="reponse">D</button>
  </div>
</section>
\end{lstlisting}

Puis on met à jour \code{afficherQuestion} pour montrer ou cacher l'image, et pour faire avancer la barre :

\begin{lstlisting}[style=js]
const elImage = document.querySelector("#image-question");
const barre   = document.querySelector("#barre-progression");

function afficherQuestion() {
  const q = questions[indexQuestion];
  elProgression.textContent = `Question ${indexQuestion + 1} / ${questions.length}`;
  elEnonce.textContent = q.enonce;

  if (q.image) {
    elImage.src = q.image;            // on met l'adresse de l'image
    elImage.classList.remove("cachee");
  } else {
    elImage.classList.add("cachee");  // pas d'image pour les autres modes
  }

  boutonsReponse.forEach((bouton, i) => {
    bouton.textContent = q.propositions[i];
  });

  // faire avancer la barre de progression
  const pourcentage = (indexQuestion / questions.length) * 100;
  barre.style.width = pourcentage + "%";
}
\end{lstlisting}

% ===================== SECTION 4 =====================
\section{La barre de progression (CSS)}

Une barre qui se remplit au fil des questions. Ajoutez à \code{style.css} :

\begin{lstlisting}[style=html]
#piste-progression {
  background: #eeeeee;
  border-radius: 6px;
  height: 10px;
  margin: 1rem 0;
}
#barre-progression {
  background: #c0392b; /* brique de la palette */
  height: 100%;
  width: 0;
  border-radius: 6px;
  transition: width 0.3s ease; /* l'animation douce du Lab 4 */
}
#image-question {
  max-width: 220px;
  display: block;
  margin: 1rem auto;
}
\end{lstlisting}

\begin{infobox}[\texttt{element.style.width}, à prendre tel quel]
En JavaScript, \code{barre.style.width = "40\%"} change \emph{directement} la largeur de l'élément. Combiné à la \code{transition} CSS, la barre se remplit en douceur. C'est le pont JS $\leftrightarrow$ CSS : le code décide, le style anime.
\end{infobox}

\begin{infobox}[Vérification]
Lancez chaque mode l'un après l'autre. En mode Drapeaux, l'image s'affiche ; dans les autres, elle disparaît. La barre avance à chaque question. \textbf{Quatre jeux pour le prix d'un moteur !}
\end{infobox}

% ===================== EXERCICES =====================
\section{Exercices d'application}

\begin{objectifbox}[Exercice 1 --- Mode « Superficie »]
Ajoutez un cinquième mode : « Quel est le plus grand de ces pays ? ».\\
\emph{Indices :} (1) ajoutez \code{area} à la liste \code{fields=} de l'URL et au \code{map} de \code{chargerPays} (champ \code{superficie}). (2) Tirez quatre pays, la bonne réponse est celui de plus grande \code{superficie}.
\end{objectifbox}

\begin{objectifbox}[Exercice 2 --- Compteur d'erreurs]
Ajoutez une variable \code{erreurs} (avec \code{let}), incrémentée dans le \code{else} de \code{repondre}, et affichez-la sur l'écran Résultat à côté du score.
\end{objectifbox}

% ===================== AIDE-MÉMOIRE =====================
\section{Aide-mémoire du Lab 9}

\begin{center}
\begin{tabular}{p{5.6cm} p{7.4cm}}
\toprule
\textbf{Besoin} & \textbf{Code} \\
\midrule
Ranger une info dans le HTML & \code{data-mode="drapeau"} \\
La relire en JS & \code{element.dataset.mode} \\
Accès par crochets & \code{objet[champ]} \\
Valeurs d'un objet & \code{Object.values(objet)} \\
Nombre $\rightarrow$ texte & \code{String(nombre)} \\
Afficher une image & \code{img.src = adresse} \\
Changer une dimension & \code{element.style.width = "40\%"} \\
\bottomrule
\end{tabular}
\end{center}

% ===================== GIT =====================
\section{Sauvegarder mon travail (routine Git)}

\begin{retenirbox}[Routine Git de fin de Lab --- tag \texttt{v0.9}]
\begin{lstlisting}[style=html]
git status
git add .
git commit -m "Lab 9 : modes Drapeaux, Langues, Voisins + progression"
git tag v0.9
git push
git push --tags
\end{lstlisting}
Vérifiez sur GitHub que le commit \emph{et} le tag \code{v0.9} apparaissent.
\end{retenirbox}

% ===================== RAPPORT =====================
\begin{attentionbox}[Rapport à produire]
Déposez \code{rapport-lab-9-NOM-Prenom.pdf} dans \code{rapports/}, avec quatre sections :
\begin{enumerate}[leftmargin=*,noitemsep]
  \item \textbf{Ce que j'ai fait} : les quatre modes, l'affichage du drapeau, la barre de progression.
  \item \textbf{Difficultés rencontrées}.
  \item \textbf{Comment je les ai résolues}.
  \item \textbf{Lien vers mon repo} avec le tag \code{v0.9}.
\end{enumerate}
\end{attentionbox}

% ===================== RÉSUMÉ =====================
\section{Résumé et transition}

Avec un seul moteur de quiz --- une fonction \code{genererQuestion} qui prend le mode en paramètre --- vous offrez quatre expériences de jeu. Vous savez désormais paramétrer un comportement, afficher des images, et animer une barre de progression.

Au \textbf{Lab 10}, on finit l'application : sauvegarder les meilleurs scores entre les sessions (\code{localStorage}), ajouter un mode chrono, et soigner l'accessibilité. Ce sera la dernière touche avant la mise en ligne.

\begin{center}
\textit{Le jeu est complet. Au Lab 10, on le rend mémorable et accessible.}
\end{center}

\bigskip
\noindent\rule{\textwidth}{0.4pt}
\begin{center}
\small Dr.~El Hadji Bassirou TOURÉ --- UCAD\\
\small Licence 1 DUT1-Informatique --- Fondamentaux du Web --- 2025--2026
\end{center}

\end{document}
